\documentclass[11pt,a4paper]{article}

\usepackage{fontspec}
\usepackage{polyglossia}
\setdefaultlanguage{russian}
\setotherlanguage{english}
\setmainfont{Arial}
\newfontfamily\cyrillicfont{Arial}
\setmonofont{Menlo}
\newfontfamily\cyrillicfonttt{Menlo}

\usepackage[a4paper,left=21mm,right=21mm,top=19mm,bottom=21mm]{geometry}
\usepackage{amsmath,amssymb}
\usepackage{booktabs}
\usepackage{graphicx}
\usepackage[section]{placeins}
\usepackage[hidelinks]{hyperref}
\usepackage{microtype}

\setlength{\parindent}{1.2cm}
\setlength{\parskip}{2pt}
\setlength{\emergencystretch}{2em}
\graphicspath{{reports/figures/}}

\title{Bias-Variance Decomposition in Practice}
\author{Исследовательский проект по регрессионным моделям}
\date{}

\begin{document}

\maketitle

\begin{abstract}
Проект исследует связь сложности регрессионной модели с bias, variance и
квадратичной ошибкой. На синтетических данных компоненты оцениваются напрямую
по независимым обучающим выборкам. На Diabetes и California Housing
используются repeated cross-validation и paired bootstrap. Для реальных данных
проверяется точное тождество квадратичной ошибки без подмены неизвестного
смещения расстоянием до эталонной модели.
\end{abstract}

\section{Синтетический эксперимент}

Пусть $Y=f(X)+\varepsilon$, где
$\mathbb E[\varepsilon\mid X]=0$ и
$\mathbb V[\varepsilon\mid X]=\sigma^2$. Для squared loss выполняется

\[
\mathbb E_{D,\varepsilon}\left[(Y-\hat f_D(x))^2\right]
=
\left(\mathbb E_D[\hat f_D(x)]-f(x)\right)^2
+
\mathbb V_D[\hat f_D(x)]
+
\sigma^2.
\]

Используется функция $f(x)=\sin(2x)+0.2x^2$ на $X\sim U[-3,3]$ и
$\sigma=0.35$. Для каждой конфигурации строятся 50 независимых обучающих
выборок по 160 объектов. Все модели получают одни и те же выборки и общую
сетку из 240 точек, которая не передается в \texttt{fit}. Случайность данных,
модели и тестового шума задается разными seed sequences. Variance Random
Forest и MLP включает изменение данных и случайность алгоритма.

На конечном числе повторов наивная оценка squared bias содержит Monte Carlo
добавку. Для $m$ повторов используется поправка

\[
\widehat{Bias^2}_{\mathrm{corr}}
=
\frac{1}{n}\sum_i\left(\bar f(x_i)-f(x_i)\right)^2
-
\frac{\widehat{Variance}}{m}.
\]

Variance оценивается с поправкой Бесселя. Из-за несмещенности исправленный
Bias$^2$ может быть немного отрицательным около нуля. Полосы на графике
получены bootstrap перевыборкой 50 prediction vectors.

\begin{figure}[htbp]
  \centering
  \includegraphics[width=\textwidth]{complexity_tradeoff.png}
  \caption{Компоненты ошибки и неопределенность Bias$^2$ и Variance}
\end{figure}

\begin{center}
\small
\begin{tabular}{llrrrr}
\toprule
Модель & Параметр & Bias$^2$ & Variance & MSE & Время, с \\
\midrule
Polynomial Ridge & degree=7 & 0.001 & 0.007 & 0.130 & 0.05 \\
MLP & hidden\_width=256 & 0.001 & 0.007 & 0.131 & 17.17 \\
KNN & n\_neighbors=12 & 0.006 & 0.015 & 0.143 & 0.05 \\
Random Forest & max\_depth=4 & 0.008 & 0.026 & 0.156 & 2.40 \\
Decision Tree & max\_depth=5 & 0.001 & 0.056 & 0.180 & 0.03 \\
\bottomrule
\end{tabular}
\end{center}


Polynomial Ridge степени 7 показывает expected MSE 0.1305. MLP ширины 256
дает 0.1312, KNN с 12 соседями дает 0.1428. У дерева corrected Bias$^2$ после
глубины 5 близок к нулю в пределах полосы, поэтому прежний визуальный хвост
нельзя интерпретировать как дальнейшее устойчивое снижение bias. Рост ошибки
дерева связан прежде всего с variance. У полиномиальной модели ошибка выходит
на плато, а у однослойной MLP не образуется U-shaped curve на исследованной
сетке.

\section{Размер выборки и выбор сложности}

Для Decision Tree глубины 6 увеличение train size с 40 до 320 объектов
снижает variance с 0.134 до 0.050. При $\sigma=0.9$ noise равен 0.81 и
становится крупнейшей частью expected MSE.

\begin{figure}[htbp]
  \centering
  \includegraphics[width=0.94\textwidth]{train_size_noise_effects.png}
  \caption{Влияние размера обучающей выборки и уровня шума}
\end{figure}

GridSearchCV теперь использует выборку того же размера, что и Monte Carlo
эксперимент. Набор уровней сложности также совпадает. CV и minimum expected
MSE выбирают одинаковые параметры для четырех семейств из пяти. Для дерева
CV выбирает глубину 4, а Monte Carlo минимум достигается при глубине 5.

\begin{center}
\small
\begin{tabular}{lrrrr}
\toprule
Модель & CV & Monte Carlo & Совпало & Время, с \\
\midrule
Polynomial Ridge & 7 & 7 & да & 0.06 \\
MLP & 256 & 256 & да & 5.66 \\
KNN & 12 & 12 & да & 0.05 \\
Random Forest & 4 & 4 & да & 1.95 \\
Decision Tree & 4 & 5 & нет & 0.03 \\
\bottomrule
\end{tabular}
\end{center}


Таблицы содержат wall-clock время полного набора fit для синтетических
конфигураций и среднее время одного CV fit для реальных данных. Эти значения
нужны для сравнения вычислительной цены, но зависят от компьютера.

\section{Нейронная сеть}

Отдельный эксперимент меняет ширину до 256, число скрытых слоев до трех и
L2 параметр $\alpha$. Для ширины от 128 до 256 expected MSE почти не меняется.
Variance также не растет. Это не подтверждает универсальный тезис о том, что
любое увеличение ширины обязательно ведет к переобучению. На этой одномерной
задаче Adam, L2 и остановка по отсутствию улучшения loss ограничивают
эффективную сложность сети. Переход от двух слоев к трем повышает variance с
0.0109 до 0.0128 и expected MSE с 0.1333 до 0.1352.

\begin{figure}[htbp]
  \centering
  \includegraphics[width=\textwidth]{mlp_capacity.png}
  \caption{Ширина, глубина и L2 как независимые оси сложности MLP}
\end{figure}

История обучения строится через \texttt{partial\_fit}. Это сохраняет
optimizer state между эпохами и не создает искусственных однопроходных
вызовов \texttt{fit}. К тысячной эпохе train RMSE снижается до 0.326, ниже
заданного $\sigma=0.35$. Validation RMSE остается около 0.374. Сеть теперь
достаточно обучена, а разрыв между кривыми показывает цену подгонки шумовой
части train sample.

\begin{figure}[htbp]
  \centering
  \includegraphics[width=0.9\textwidth]{mlp_training_history.png}
  \caption{Train и validation RMSE при последовательном обучении MLP}
\end{figure}

\section{Реальные данные}

Diabetes содержит 442 наблюдения и 10 признаков. California Housing содержит
20640 наблюдений. Для практичного полного запуска один раз выбираются 5000
строк по фиксированному seed. Исходный датасет загружается средствами
scikit-learn и остается в локальном cache. Масштабирование Ridge и MLP
находится внутри Pipeline и обучается заново на каждом fold.

Модели сравниваются на одинаковых разбиениях RepeatedKFold с 5 folds и 3
повторами. На Diabetes Ridge показывает RMSE $54.86\pm3.05$. MLP и Gradient
Boosting находятся в пределах наблюдаемого разброса. На California Housing
MLP показывает RMSE $0.570\pm0.019$, Random Forest дает $0.612\pm0.021$.
Dummy заметно хуже обучаемых моделей на обоих наборах.

\begin{figure}[htbp]
  \centering
  \includegraphics[width=\textwidth]{real_model_comparison.png}
  \caption{Repeated cross-validation на двух реальных наборах}
\end{figure}

\begin{center}
\small
\begin{tabular}{llrrrr}
\toprule
Датасет & Модель & RMSE & Std & $R^2$ & Время, с \\
\midrule
California Housing & MLP & 0.570 & 0.019 & 0.758 & 0.470 \\
California Housing & Random Forest & 0.612 & 0.021 & 0.721 & 0.230 \\
California Housing & Gradient Boosting & 0.696 & 0.015 & 0.640 & 0.416 \\
California Housing & Ridge & 0.759 & 0.041 & 0.570 & 0.001 \\
California Housing & Decision Tree & 0.763 & 0.019 & 0.568 & 0.007 \\
California Housing & Dummy & 1.161 & 0.022 & -0.001 & 0.000 \\
Diabetes & Ridge & 54.857 & 3.046 & 0.486 & 0.001 \\
Diabetes & MLP & 55.982 & 4.321 & 0.464 & 0.028 \\
Diabetes & Gradient Boosting & 56.533 & 3.163 & 0.455 & 0.045 \\
Diabetes & Random Forest & 57.349 & 3.916 & 0.438 & 0.080 \\
Diabetes & Decision Tree & 64.303 & 3.789 & 0.293 & 0.001 \\
Diabetes & Dummy & 77.159 & 3.155 & -0.016 & 0.001 \\
\bottomrule
\end{tabular}
\end{center}


\section{Сложность на Diabetes}

Для одного train-test split используются общие 30 bootstrap samples. У Ridge
усиление регуляризации снижает variance со 106 до 51, при этом средний
bootstrap MSE остается около 2885. У дерева минимум среднего bootstrap MSE
наблюдается при глубине 2, после чего variance быстро растет. Для KNN variance
падает с ростом числа соседей. У MLP ширина 64 устойчивее ширины 128 при всех
трех значениях $\alpha$. Таким образом, реальная часть исследует не только
фиксированные конфигурации, но и траектории сложности на десяти признаках.

\begin{figure}[htbp]
  \centering
  \includegraphics[width=\textwidth]{diabetes_complexity.png}
  \caption{Ошибка и bootstrap variance при изменении сложности на Diabetes}
\end{figure}

\section{Точное bootstrap тождество}

Истинная regression function на реальных данных неизвестна. Поэтому bias и
irreducible noise раздельно не оцениваются. Для фиксированного test set и
prediction matrix проверяется точное тождество

\[
\frac{1}{B}\sum_{b=1}^{B}MSE(y,\hat f_b)
=
MSE(y,\bar f)
+
\frac{1}{n}\sum_i Var_b\left(\hat f_b(x_i)\right).
\]

Первое слагаемое объединяет систематическую ошибку относительно наблюдаемых
targets, test noise и эффект выбранного split. Его нельзя называть Bias$^2$.
Второе слагаемое измеряет нестабильность прогноза при перевыборке. Identity
gap во всех строках находится на уровне машинной точности.

\begin{figure}[htbp]
  \centering
  \includegraphics[width=\textwidth]{real_bootstrap_identity.png}
  \caption{Точное тождество для paired bootstrap predictions}
\end{figure}

\section{Выводы и ограничения}

Классический bias-variance механизм ясно виден у Decision Tree, KNN и Random
Forest. Для Polynomial Ridge после достаточной степени возникает плато. У MLP
ширина сама по себе не вызывает роста variance на синтетической задаче, но
дополнительная глубина и реальный многомерный набор показывают менее устойчивые
режимы. Поэтому общий вывод формулируется по наблюдаемым траекториям, а не как
универсальное свойство всех пяти семейств.

Оценки зависят от конечного числа Monte Carlo и bootstrap повторов. Полосы
неопределенности отражают перевыборку prediction vectors, а не полный повтор
всего исследования. Время обучения зависит от оборудования. На реальных
данных тождество разделяет ошибку среднего прогноза и нестабильность, но не
разделяет bias и noise без дополнительных предположений о неизвестной
регрессионной функции.

\end{document}
